\section{Pattern 1：Tool Wrapper（工具封装）}

Tool Wrapper 为 Agent 提供针对特定库的\textbf{按需上下文}。与其将 API 约定硬编码到 system prompt 中，不如将它们打包成一个 skill，Agent 只在实际使用该技术时才加载相关上下文。

\begin{importantbox}{Tool Wrapper 的核心机制}
\texttt{SKILL.md} 文件监听用户 prompt 中的特定库关键词，动态从 \texttt{references/} 目录加载内部文档，并将这些规则作为绝对真理应用。这正是将团队内部编码规范或特定框架最佳实践直接分发到开发者工作流中的机制。
\end{importantbox}

\subsection{工作原理}

Tool Wrapper 是最简单的模式。其核心流程为：

\begin{enumerate}
    \item Agent 检测到用户 prompt 中包含特定库的关键词（如 ``FastAPI''）
    \item \texttt{SKILL.md} 指令触发，动态加载 \texttt{references/} 目录下的约定文档（如 \texttt{conventions.md}）
    \item Agent 将加载的规则作为上下文，指导代码编写或审查
\end{enumerate}

\subsection{适用场景}

\begin{itemize}
    \item 分发团队内部的编码规范
    \item 封装特定框架的最佳实践
    \item 在不同项目间复用一致的技术标准
\end{itemize}

\begin{knowledgebox}{按需加载的优势}
Tool Wrapper 不会在所有对话中都占用 context window。只有当用户的请求涉及特定技术栈时，相关文档才会被加载。这种按需加载策略既节省了 token 开销，又保证了上下文的精准性。
\end{knowledgebox}

\subsection{本章小结}

Tool Wrapper 是最简单也最基础的模式，通过关键词触发和按需加载，将特定库的知识动态注入 Agent 的上下文中，适合封装编码规范和框架约定。

